\begin{table}[H]
\centering\footnotesize
\setlength{\tabcolsep}{4pt}
\caption{Hungria: tamanho e financiamento dos governos locais (setor S1313)}
\label{tab:hungria}
\begin{adjustbox}{max width=\textwidth}
\begin{tabular}{lrrrrrrr}
\toprule
Indicador & 2012 & 2013 & 2014 & 2016 & 2019 & 2022 & 2025 \\
\midrule
Despesa dos governos locais (\% do PIB) & 9,2 & 7,5 & 7,8 & 5,9 & 6,7 & 5,9 & 5,4 \\
Receita dos governos locais (\% do PIB) & 9,7 & 10,0 & 9,1 & 6,2 & 6,6 & 5,7 & 5,5 \\
Despesa local / desp. do gov. geral & 18,7\% & 15,0\% & 15,6\% & 12,7\% & 14,7\% & 12,1\% & 11,3\% \\
Transf. correntes recebidas (\% do PIB) & 4,6 & 3,0 & 2,9 & 2,5 & 2,1 & 2,1 & 2,0 \\
Transf. de capital recebidas (\% do PIB) & 1,3 & 3,4 & 3,0 & 0,6 & 1,3 & 1,0 & 0,6 \\
Impostos locais (\% do PIB) & 2,4 & 2,2 & 2,1 & 2,2 & 2,1 & 1,7 & 2,0 \\
Transf. correntes / receita sem transf. de capital & 55,5\% & 46,1\% & 47,1\% & 44,6\% & 39,9\% & 44,2\% & 41,2\% \\
Transf. correntes e de capital / receita & 61,7\% & 64,6\% & 64,7\% & 49,5\% & 52,1\% & 54,1\% & 47,1\% \\
Despesa local per capita (mil HUF de 2025) & 475 & 401 & 452 & 383 & 529 & 517 & 492 \\
\bottomrule
\end{tabular}
\end{adjustbox}
\par\vspace{3pt}\parbox{\textwidth}{\scriptsize Nota: governos locais = setor S1313 das contas nacionais (SEC 2010). Transferências correntes = D73 (dentro do governo geral); de capital = D9 (em 2013--2014 inclui a assunção das dívidas municipais pelo governo central); impostos = D2 + D5 + D91. Período: 2012--2025; a série obtida começa já dentro da primeira onda de recentralização. Valores reais deflacionados pelo HICP médio anual (preços de 2025) e divididos pela população (Banco Mundial). Fonte: elaboração própria com dados do Eurostat (gov\_10a\_main, prc\_hicp\_aind) e do Banco Mundial.}
\end{table}
